\section{Conclusiones}

En este trabajo se estudió la aplicación del algoritmo L-BFGS al problema de planificación inversa en radioterapia IMRT, analizando tanto sus fundamentos matemáticos como su implementación en casos clínicos reales. A partir de esto, fue posible relacionar conceptos desarrollados en métodos numéricos con un problema concreto de física médica que presenta una complejidad computacional muy elevada.

Uno de los aspectos más importantes del trabajo fue entender por qué los métodos quasi-Newton resultan especialmente útiles en este tipo de aplicaciones. El problema de planificación inversa involucra miles de variables asociadas a las intensidades de irradiación, por lo que métodos que requieren almacenar o invertir el Hessiano completo se vuelven impracticables. En este contexto, L-BFGS logra incorporar información de curvatura manteniendo un costo computacional reducido, lo que explica su amplio uso en optimización de gran escala.

También resultó interesante observar cómo un problema físico y clínico puede formularse matemáticamente mediante funciones objetivo que representan distintos criterios dosimétricos. La cobertura tumoral, la protección de órganos de riesgo y la reducción de dosis sobre tejido sano son objetivos que compiten entre sí, por lo que no existe una única solución óptima. El análisis de los frentes de Pareto mostró justamente este compromiso entre objetivos, y cómo distintas soluciones pueden ser igualmente válidas dependiendo del criterio clínico adoptado.

Los resultados obtenidos muestran que L-BFGS presenta una convergencia muy robusta para los casos analizados. En prácticamente todas las ejecuciones, el algoritmo convergió al mismo resultado aun utilizando condiciones iniciales aleatorias diferentes. Además, la comparación con Fast Simulated Annealing permitió verificar que las soluciones obtenidas corresponden efectivamente al mínimo global del problema, pero con tiempos de cálculo muchísimo menores.

Otro aspecto importante fue observar que la optimización matemática no reemplaza el criterio clínico, sino que funciona como una herramienta de apoyo para la toma de decisiones. Los histogramas dosis-volumen y los frentes de Pareto muestran que siempre existe un compromiso entre irradiar adecuadamente el tumor y minimizar el daño sobre estructuras sanas. En consecuencia, la elección final del tratamiento sigue dependiendo de prioridades médicas específicas para cada paciente.

En conjunto, el trabajo permitió mostrar cómo herramientas clásicas de optimización numérica encuentran aplicaciones directas en tecnologías médicas modernas de alta complejidad. Además de estudiar el funcionamiento del algoritmo L-BFGS, el desarrollo permitió entender mejor la relación entre modelado matemático, eficiencia computacional y aplicaciones reales dentro de la física médica.